% Source PNG: img/2007/guangdong_science/q6_flowchart.png; source vector/text
% flowchart from source/普通高考/2007/2007广东理.pdf, page 1, question 6.
% Node -> directed-edge list: 开始 -> 输入 A_1,...,A_{10} -> s=0,i=4 -> blank test;
% test(是) -> s=s+A_i -> i=i+1 -> left return arrow to the stem above the test;
% test(否) -> 输出s -> 结束. The condition is intentionally blank in the source.

\begin{tikzpicture}[
  >=Stealth,
  x=1cm,
  y=0.92cm,
  line width=0.45pt,
  line cap=round,
  line join=round,
  font=\normalsize,
  flowtext/.style={anchor=center, align=center,
    execute at begin node={\everypar{}\leftskip=0pt\rightskip=0pt\parindent=0pt\hangindent=0pt\hangafter=1}},
  every node/.style={flowtext},
  flowbox/.style={draw, flowtext, minimum height=0.72cm,
    inner xsep=3pt, inner ysep=2pt},
  terminal/.style={flowbox, rounded corners=0.18cm, minimum width=1.45cm},
  io/.style={flowbox, trapezium, trapezium left angle=72,
    trapezium right angle=108, minimum height=0.80cm},
  process/.style={flowbox},
  decision/.style={flowbox, diamond, aspect=2.8, minimum width=3.70cm,
    minimum height=1.06cm, inner sep=0pt},
  caption/.style={flowtext, inner sep=0pt},
]
  \node[terminal] (start) at (0,13.15) {开始};
  \node[io, minimum width=6.35cm] (input) at (0,11.65)
    {输入 $A_1,A_2,\ldots,A_{10}$};
  \node[process, minimum width=3.30cm] (init) at (0,10.05)
    {$s=0,i=4$};
  \node[decision] (test) at (0,7.45) {};
  \node[process, minimum width=3.20cm] (add) at (3.55,7.45)
    {$s=s+A_i$};
  \node[process, minimum width=2.75cm] (increment) at (3.55,8.75)
    {$i=i+1$};
  \node[io, minimum width=2.15cm] (output) at (0,4.95) {输出 $s$};
  \node[terminal] (finish) at (0,3.35) {结束};

  \coordinate (loop-join) at (0,8.75);
  % Keep the return arrowhead 3 mm before the shared stem junction; the
  % final unarrowed connector preserves the original loop topology.
  \coordinate (return-tip) at (0.35,8.75);
  \coordinate (right-turn) at (3.55,7.45);

  \draw[->] (start.south) -- (input.north);
  \draw[->] (input.south) -- (init.north);
  \draw (init.south) -- (loop-join);
  \draw[->] (loop-join) -- (test.north);

  % 是 branch: update s, increment i, then return left into the main stem.
  \draw[->] (test.east) -- (right-turn) -- (add.west);
  \draw[->] (add.north) -- (increment.south);
  \draw[->] (increment.west) -- (return-tip);
  \draw (return-tip) -- (loop-join);

  % 否 branch terminates independently; no extra arrow on the horizontal return.
  \draw[->] (test.south) -- node[flowtext, right=2pt] {否} (output.north);
  \draw[->] (output.south) -- (finish.north);
  \node[flowtext, anchor=west] at (1.75,7.95) {是};
  \node[caption] at (0,2.25) {图2};
\end{tikzpicture}
